\chapter{Introduction}
\label{ch:introduction}

Modern software systems are large, rapidly evolving, and security-critical.
Static Application Security Testing (SAST) tools analyse source code without
executing it and can detect a wide range of security vulnerabilities early in
the development lifecycle \cite{hegedus2022sca}. Tools such as CodeQL, Semgrep,
and SonarQube have become first-class citizens in DevSecOps workflows, integrated
directly into CI/CD pipelines to catch vulnerabilities before deployment.

Despite their many advantages, SAST tools suffer from a well-documented and
persistent limitation: they produce a disproportionately large number of
\emph{false positive alerts}. The ratio of false positive warnings can reach
30--60\% \cite{hegedus2022sca}. This high false-alarm rate has serious
consequences. Developers become overwhelmed and begin to distrust the tools,
eventually ignoring most warnings or disabling the tools entirely --- a
phenomenon known as \emph{alert fatigue}. Real vulnerabilities sink to the
bottom of unvisited alert queues and remain unresolved, defeating the purpose
of running static analysis.

The core challenge is that the same code pattern may be a genuine vulnerability
in one context and a benign construct in another. Classical rule-based SAST tools
cannot capture context-dependent differences because their detection logic is
fixed and shallow. What is needed is a post-processing layer that can learn from
historical developer decisions to distinguish actionable warnings from false alarms.

This project proposes and evaluates an \textbf{Agentic AI Pipeline for Automated
Vulnerability Detection and Triage} with three complementary layers:

\begin{enumerate}
  \item \textbf{Detection layer}: CodeQL SAST analysis on OWASP Benchmark Java
        test cases, producing SARIF-format results.
  \item \textbf{Triage layer}: Machine Learning classifiers trained on features
        extracted from SARIF alerts and source code context, classifying each
        alert as a true or false positive and assigning a confidence score.
  \item \textbf{Agentic layer}: An agentic AI component that consumes
        high-confidence alerts, generates remediation suggestions, and supports
        automated patch generation and re-evaluation.
\end{enumerate}

This approach directly addresses the gaps identified in prior work by incorporating 
a hybrid system combining static analysis signals with ML-based filtering and Large 
Language Model (LLM) reasoning for patching.

The remainder of this report is structured as follows:
Chapter \ref{ch:literature} reviews the literature.
Chapter \ref{ch:methods} details the system methodology, including threshold optimization.
Chapter \ref{ch:results} presents experimental results and model comparisons.
Chapter \ref{ch:conclusion} concludes with future directions.
